% GENERATED_BY: oc_core_monograph_translation_draft_pipeline_v1.py
% AI_ASSISTED_TRANSLATION_DRAFT: TRUE
% THEOREM_REVIEW_STATUS: PENDING
% THEOREM_REVIEW_MODE: FORMAL_AI_GATE__CODEX_CLI_CHATGPT
% THEOREM_REVIEW_REVIEWER_ID:
% THEOREM_REVIEW_AUTHORITY_CLASS:
% THEOREM_REVIEWED_AT:
% THEOREM_REVIEW_DOSSIER_REF:
% SOURCE_LANGUAGE: EN
% TARGET_LANGUAGE: RU
% SOURCE_EN_REF: content/m_spaces/m4.tex
% ================================================================
% ==== FILE: content/m_spaces/m4.tex
% ================================================================

% ==============================
%  Ontology of Continua — Core
%  Meta-Space: M4
%  FULL MODULE — FINAL VERSION
% ==============================

\subsubsection{\texorpdfstring{$M_4$}{M_4} Обзор}
\label{sec:m4-overview}

$M_4$ — это мета-пространство, задающее условия допустимости для
биологических континуумов уровня $K_4$.
Это первое мета-пространство, в котором:

\begin{itemize}
    \item допустима стабильная \emph{полупроницаемая граница},
    \item топологическое различие \textbf{внутри/снаружи} становится осью,
    \item могут поддерживаться не-тривиальные \\emph{градиенты концентрации, электрического,
    pH и редокс потенциалов},
    \item разрешены процессы активного транспорта и преобразования энергии,
    \item реакционные сети локализуются и пространственно координируются,
    \item переход от химической к прото-биологической структуре ($K_3 \to K_4$)
          допустим.
\end{itemize}

$M_4$ вводит геометрические и энергетические условия
компартментации — определяющего признака ранних клеток.

% ---------------------------------------------------------------
\subsubsection*{1. Допустимое конфигурационное пространство \texorpdfstring{$\Omega(M_4)$}{\Omega(M_4)}}

Допустимые конфигурации в $M_4$ расширяют конфигурации $M_3$ добавлением:

\[
\Omega(M_4)
=
\Big\{
(\phi,\ G_{\mathrm{chem}},\ \rho_{\mathrm{in/out}},
\Delta\mu,\ g_{ij},\ \partial\Omega,\ \Pi)
\ \Big|\ \text{выполняются условия допустимости}
\Big\}.
\]

Где:

\begin{itemize}
    \item $\partial\Omega$ — полупроницаемая мембрана, формирующая компактную границу,
    \item $\rho_{\mathrm{in/out}}$ — концентрации внутри и снаружи,
    \item $\Delta\mu$ — градиенты электрохимического потенциала через $\partial\Omega$,
    \item $\Pi$ — тензор проницаемости (направленно-зависимый и зависимый от вида частиц),
    \item $G_{\mathrm{chem}}$ — внутренние RAF-сети, остающиеся функциональными
          под конфайнментом.
\end{itemize}

Условия допустимости:

\begin{enumerate}
    \item \textbf{Стабильность мембраны}:
    $\partial\Omega$ должен сохранять механическую целостность при
    натяжении
    $T_{\mathrm{mem}} < \Theta_{\mathrm{mem}}$.

    \item \textbf{Ограничения проницаемости}:
    $\Pi$ должен быть конечным, анизотропным и видово-специфичным.

    \item \textbf{Согласованность градиентов}:
    электрохимические градиенты должны удовлетворять:
    \[
    |\Delta\mu| < \Theta_{\mathrm{grad-max}}.
    \]

    \item \textbf{Конечность объёма}:
    внутренний регион должен иметь конечный объём и массу.

    \item \textbf{RAF-совместимость}:
    внутренние реакционные сети остаются топологически замкнутыми.
\end{enumerate}

$M_4$ поэтому задаёт минимальную среду, в которой могут существовать
протоклетки.

% ---------------------------------------------------------------
\subsubsection*{2. Допустимые оси \texorpdfstring{$A(M_4)$}{A(M_4)}}

Появляется новая базовая ось:

\[
A_{\mathrm{in/out}} = \{ \text{inside},\ \text{outside} \}.
\]

Она удовлетворяет:

\[
A_{\mathrm{in/out}} \not\subseteq \operatorname{span}(A_1, A_2, A_3)
\quad\Rightarrow\quad
\text{допускается новая размерность}.
\]

К этой группе относятся также:

\begin{itemize}
    \item $A_{\mathrm{perm}}$ — состояния проницаемости,
    \item $A_{\mathrm{grad}}$ — конфигурации градиентов (pH, ионы, заряд),
    \item $A_{\mathrm{vol}}$ — ось объёма/давления,
    \item $A_{\mathrm{energy}}$ — внутренние энергетические состояния
          (предшественники ATP, редокс-пулы).
\end{itemize}

Эти оси определяют внутреннюю функциональную геометрию ранних клеток.

% ---------------------------------------------------------------
\subsubsection*{3. Потенциалы \texorpdfstring{$P(M_4)$}{P(M_4)}}

$M_4$ допускает новые классы потенциалов:

\[
P(M_4) =
\big(
U_{\mathrm{mem}},\
U_{\mathrm{grad}},\
U_{\mathrm{osm}},\
U_{\mathrm{charge}},\
U_{\mathrm{redox}},\
U_{\mathrm{active}}
\big).
\]

Интерпретация:

\begin{itemize}
    \item $U_{\mathrm{mem}}$ — потенциал изгиба мембраны и поверхностного натяжения,
    \item $U_{\mathrm{grad}}$ — энергия, запасённая в химических и электрических градиентах,
    \item $U_{\mathrm{osm}}$ — потенциал осмотического давления,
    \item $U_{\mathrm{charge}}$ — кулоновские и диэлектрические взаимодействия,
    \item $U_{\mathrm{redox}}$ — редокс-энергетический ландшафт,
    \item $U_{\mathrm{active}}$ — преобразование энергии, обеспечивающее
          активный транспорт.
\end{itemize}

Все потенциалы должны быть:

\begin{enumerate}
    \item почти всюду дифференцируемыми,
    \item ограниченными снизу,
    \item совместимыми с условиями стабильности мембраны.
\end{enumerate}

% ---------------------------------------------------------------
\subsubsection*{4. Пороги \texorpdfstring{$\Theta(M_4)$}{\Theta(M_4)}}

Ключевые биологические пороги:

\begin{itemize}
    \item $\Theta_{\mathrm{mem}}$ — порог разрыва мембраны,
    \item $\Theta_{\mathrm{grad}}$ — минимальный градиент для поддержания структуры,
    \item $\Theta_{\mathrm{grad-max}}$ — максимально допустимый градиент,
    \item $\Theta_{\mathrm{osm}}$ — осмотическая толерантность,
    \item $\Theta_{\mathrm{pH}}$ — интервал жизнеспособности внутреннего pH,
    \item $\Theta_{\mathrm{redox}}$ — порог редокс-совместимости.
\end{itemize}

Критическое условие сохранения живучести $K_4$:

\[
T_4 <
\min\{
\Theta_{\mathrm{mem}},\
\Theta_{\mathrm{osm}},\
\Theta_{\mathrm{grad-max}},\
\Theta_{\mathrm{pH}},\
\Theta_{\mathrm{redox}}
\}.
\]

% ---------------------------------------------------------------
\subsubsection*{5. Потоки \texorpdfstring{$J(M_4)$}{J(M_4)}}

Допустимые потоки:

\[
J(M_4)
=
\big(
J_{\mathrm{diff}},\
J_{\mathrm{perm}},\
J_{\mathrm{pump}},\
J_{\mathrm{redox}},\
J_{\mathrm{osm}},\
\nabla_i T_4
\big).
\]

Где:

\begin{itemize}
    \item $J_{\mathrm{diff}}$ — пассивная диффузия внутри/снаружи,
    \item $J_{\mathrm{perm}}$ — мембранные пермеационные потоки,
    \item $J_{\mathrm{pump}}$ — активный перенос против градиентов,
    \item $J_{\mathrm{redox}}$ — редокс-циклические потоки,
    \item $J_{\mathrm{osm}}$ — осмотические балансные потоки,
    \item $\nabla_i T_4$ — потоки по градиенту структурной натяжки.
\end{itemize}

Все потоки должны удовлетворять сохранению массы внутри границы мембраны.

% ---------------------------------------------------------------
\subsubsection*{6. Циклы в \texorpdfstring{$M_4$}{M_4}}

$M_4$ — это мета-пространство, где становятся допустимыми биологические циклы:

\begin{itemize}
    \item $C_{\mathrm{RAF-core}}$ — реакционно-автокаталитические циклы,
    \item $C_{\mathrm{metabolic}}$ — базовый метаболический оборот,
    \item $C_{\mathrm{pump}}$ — циклы активного транспорта насосов,
    \item $C_{\mathrm{redox}}$ — циклы переноса электронов,
    \item $C_{\mathrm{buffer}}$ — pH-буферные циклы.
\end{itemize}

Условие устойчивости цикла:

\[
\oint_C dA_i \approx 0
\quad\text{и}\quad
T_4 \ \text{ограничена}.
\]

Эти циклы формируют временную структуру $\tau(K_4)$.

% ---------------------------------------------------------------
\subsubsection*{7. Граница \texorpdfstring{$\partial\Omega(M_4)$}{\partial\Omega(M_4)}}

Конфигурация приближается к $\partial\Omega(M_4)$, когда:

\begin{itemize}
    \item мембранное натяжение приближается к порогу разрыва:
    \[
    T_{\mathrm{mem}} \to \Theta_{\mathrm{mem}},
    \]
    \item внутренний pH выходит за пределы жизнеспособного интервала,
    \item градиенты превосходят $\Theta_{\mathrm{grad-max}}$,
    \item осмотический дисбаланс становится неустойчивым,
    \item конфайнмент нарушает RAF-замыкание внутри компартмента.
\end{itemize}

Пересечение $\partial\Omega(M_4)$ разрушает протобиологический континуум $K_4$.

% ---------------------------------------------------------------
\subsubsection*{8. Связь с \texorpdfstring{$M_3$}{M_3} и \texorpdfstring{$M_5$}{M_5}}

Переход $K_3 \to K_4$ допустим, когда:

\[
A_{\mathrm{in/out}} \in A(M_4),
\qquad
\Omega(K_4) \subseteq \Omega(M_4),
\qquad
T_3 \ge \Theta_{\mathrm{dim}}.
\]

Переход к $K_5$ требует:

\[
A_{\mathrm{exc}} \in A(M_5),
\qquad
\Delta V \ge \Theta_{\mathrm{exc}},
\qquad
\text{допустимы каналы и электрическая проводимость}.
\]

Итак, $M_4$ — мета-пространство:

\begin{itemize}
    \item структуры мембраны,
    \item компартментации,
    \item градиентов и осмотической физики,
    \item преобразования энергии,
    \item раннего метаболизма,
    \item минимальной биологической устойчивости.
\end{itemize}

Это структурная основа, обеспечивающая существование протоклеток.
